\uuid{iWdQ}
\titre{Équation homogène/solution particulière/condition initiale}
\auteur{hanani ; mzari}
\organisation{exo7}
\licence{CC BY-NC-SA 4.0}
\datecreate{2020-01-20}
\niveau{L1}
\module{Analyse}
\chapitre{Equation différentielle}
\sousChapitre{Résolution d'équation différentielle du premier ordre}
\difficulte{2}
\theme{Equations du premier ordre, équation homogène, solution particulière, condition initiale}
\exo7qcmid{lille-2.equations-differentielles.005}

\contenu{
\begin{question}
\qclassification{225.01}
\qtags{motcle=équation homogène/solution particulière/condition initiale}
On considère l'équation différentielle $(E)$ : $\displaystyle y'+y=\mathrm{e}^{x}$. Parmi les affirmations suivantes cocher celles qui sont vraies :
\begin{answers}
\bad{La solution générale de l'équation homogène est $y=k\mathrm{e}^{x}$, $k\in \Rr$.}
\bad{La fonction $\displaystyle y_0=\mathrm{e}^{x}$ est une solution particulière de $(E)$.}
\good{La solution de $(E)$ vérifiant $y(0)=0$ est $\displaystyle y=\frac{\mathrm{e}^{x}-\mathrm{e}^{-x}}{2}$.}
\good{La solution de $(E)$ vérifiant $y'(0)=0$ est $\displaystyle y=\frac{\mathrm{e}^{x}+\mathrm{e}^{-x}}{2}$.}
\end{answers}
\begin{explanations}
La solution générale de l'équation homogène est $y=k\mathrm{e}^{x}$, $k\in \Rr$. On cherche une solution particulière sous la forme $\displaystyle y_0=a\mathrm{e}^{x}$. La solution générale de $(E)$ est $\displaystyle y=k\mathrm{e}^{-x}+\frac{\mathrm{e}^{x}}{2}$, $k\in \Rr$. La condition $y(0)=0$ donne $k=-1/2$, et la condition $y'(0)=0$ donne $k=1/2$.
\end{explanations}
\end{question}
}
